% ============================================================
% Model: 自由电子气与 Drude 输运
% ============================================================

\section{自由电子气与 Drude 输运}
\label{sec:free-electron-gas}

\begin{tipbox}[关键词标签]
费米能 · 费米波矢 · Drude 模型 · 弛豫时间 · 漂移速度 · 平均自由程 · 能态密度
\end{tipbox}

% --------------------------------------------------------
% 1. 模型定义框
% --------------------------------------------------------
\begin{modelbox}[模型：自由电子气]
\textbf{物理情境}：金属中价电子脱离原子核束缚，在正离子背景中自由运动。忽略电子--电子相互作用和电子--离子周期性势场（或将其平均掉），电子视为自由粒子。

\textbf{基本假设}（Sommerfeld 模型 + Drude 近似）：
\begin{itemize}[nosep]
    \item 自由电子近似：忽略周期势，$E(\mathbf{k})=\hbar^2k^2/2m$
    \item 独立电子近似：忽略电子--电子相互作用
    \item 弛豫时间近似：碰撞使分布以指数形式恢复平衡，特征时间 $\tau$
    \item 泡利不相容：电子服从费米--狄拉克统计
\end{itemize}

\textbf{关键参数}：
\begin{itemize}[nosep]
    \item $n$：电子浓度（由原子密度和价电子数决定）
    \item $E_F, k_F, v_F$：费米能、费米波矢、费米速度
    \item $\tau, l, v_d$：弛豫时间、平均自由程、漂移速度
    \item $\rho, \sigma$：电阻率、电导率
\end{itemize}
\end{modelbox}

% --------------------------------------------------------
% 2. 关键公式框
% --------------------------------------------------------
\begin{formulabox}[核心公式]
\begin{enumerate}[label=\textbf{(\arabic*)}]
    \item \textbf{费米参量}：
    \[
    k_F=(3\pi^2 n)^{1/3},\quad
    E_F=\frac{\hbar^2k_F^2}{2m}=\frac{\hbar^2}{2m}(3\pi^2 n)^{2/3},\quad
    v_F=\frac{\hbar k_F}{m}=\sqrt{\frac{2E_F}{m}}
    \]

    \item \textbf{Drude 输运公式}：
    \[
    \sigma=\frac{ne^2\tau}{m}=\frac{1}{\rho},\qquad
    \tau=\frac{m}{ne^2\rho}=\frac{\mu m}{e}
    \]

    \item \textbf{漂移速度}：
    \[
    v_d=\frac{e\mathcal{E}\tau}{m}=\frac{\mathcal{E}}{ne\rho}=\mu\mathcal{E}
    \]
    \texteq{$\mu=e\tau/m$ 为迁移率}

    \item \textbf{平均自由程}：
    \[
    l=v_F\tau=\frac{m v_F}{ne^2\rho}
    \]
    \texteq{电子以费米速度运动，两次碰撞间走过的距离}

    \item \textbf{各维度能态密度}：
    \[
    g_{3D}(E)=\frac{V}{2\pi^2}\Bigl(\frac{2m}{\hbar^2}\Bigr)^{3/2}\sqrt{E},\;
    g_{2D}(E)=\frac{Am}{\pi\hbar^2},\;
    g_{1D}(E)=\frac{L}{\pi}\sqrt{\frac{2m}{\hbar^2}}\frac{1}{\sqrt{E}}
    \]
\end{enumerate}
\end{formulabox}

% --------------------------------------------------------
% 3. 训练点框
% --------------------------------------------------------
\begin{drillbox}[训练点与考法变体]
\trainingpoint{1} \textbf{费米参量链式计算}：给定材料密度、原子量、价电子数，计算 $n\to k_F\to E_F\to v_F$，并估算 $E_F/k_BT$（判断简并程度）。

\trainingpoint{2} \textbf{Drude 公式链}：给定 $\rho, n$，求 $\tau\to v_d\to l$；或给定 $\tau$，求 $\rho$ 和 $\sigma$。

\trainingpoint{3} \textbf{能态密度应用}：
\begin{itemize}[nosep]
    \item 3D：$N(E)=\int_0^E g_{3D}(E')dE'\propto E^{3/2}$
    \item 2D：电子数 $N=g_{2D}\cdot E$（线性关系，量子霍尔效应基础）
    \item 1D：范霍夫奇点 $g_{1D}(E)\propto E^{-1/2}$
\end{itemize}

\trainingpoint{4} \textbf{与实验对比}：由铜的实测电阻率估算 $\tau$ 和 $l$，判断室温下 $l\sim$ 原子间距（验证 Drude 模型的合理性）。

\trainingpoint{5} \textbf{磁场中的自由电子}：朗道能级 $E_n=\hbar\omega_c(n+1/2)$，磁化率振荡（de Haas--van Alphen 效应），回旋共振 $\omega_c=eB/m^*$。
\end{drillbox}

% --------------------------------------------------------
% 4. 解题技巧框
% --------------------------------------------------------
\begin{tipbox}[快速识别与解题技巧]
\begin{itemize}
    \item \examnote{识别标志}：题干出现``铜''''钠''''金属''''自由电子''''费米能''''电阻率''''漂移速度''，立即定位本模型。
    \item \examnote{费米参量速查}：
    \begin{center}
    \begin{tabular}{ll}
    \toprule
    \textbf{量} & \textbf{公式} \\
    \midrule
    $k_F$ & $(3\pi^2 n)^{1/3}$ \\
    $E_F$ & $\hbar^2k_F^2/2m$ \\
    $v_F$ & $\hbar k_F/m$ \\
    $\lambda_F$ & $2\pi/k_F$ \\
    \bottomrule
    \end{tabular}
    \end{center}
    \item \examnote{Drude 链记忆}：$\rho\xrightarrow{\tau=m/ne^2\rho}\tau\xrightarrow{v_d=e\mathcal{E}\tau/m}v_d\xrightarrow{l=v_F\tau}l$。
    \item \examnote{量纲速算}：
    \begin{itemize}[nosep]
        \item $[x]=L$：$k_F^{-1}\sim$ 费米波长 $\lambda_F/2\pi$
        \item $[E]$：$E_F$ 典型值 $1\sim10\,\text{eV}$（金属），远大于 $k_BT\sim0.025\,\text{eV}$
        \item $[\tau]$：$10^{-14}\sim10^{-13}\,\text{s}$（室温金属）
        \item $[l]$：$1\sim10\,\text{nm}$（室温），低温可达 $\mu\text{m}\sim\text{mm}$
    \end{itemize}
    \item \examnote{简并判断}：$E_F\gg k_BT$ 时电子气强简并，只有费米面附近 $k_BT$ 范围内的电子参与输运。
\end{itemize}
\end{tipbox}

% --------------------------------------------------------
% 5. 易错点框
% --------------------------------------------------------
\begin{errorbox}{常见失误与陷阱}
\begin{itemize}
    \item \textbf{单位换算}：电阻率常给 $\Omega\cdot\text{cm}$，必须换算为 $\Omega\cdot\text{m}$（乘 $10^{-2}$）。浓度常给 $\text{cm}^{-3}$，必须换算为 $\text{m}^{-3}$（乘 $10^6$）。
    \item \textbf{弛豫时间公式}：$\tau=m/(ne^2\rho)$，不要漏掉电子质量 $m$ 或误写为 $m^*$（自由电子模型中用裸质量 $m$）。
    \item \textbf{漂移速度 vs 费米速度}：$v_d\sim 10^{-4}v_F$，电场只使费米球整体微小偏移。不要把 $v_d$ 和 $v_F$ 混淆。
    \item \textbf{平均自由程}：$l=v_F\tau$（用费米速度），不是 $l=v_d\tau$。
    \item \textbf{能态密度单位}：$g(E)$ 是单位能量的态数，注意区分 ``单位体积态密度'' $g(E)/V$ 与总态数 $N(E)$。
    \item \textbf{2D 态密度}：$g_{2D}(E)=Am/(\pi\hbar^2)$ 与能量无关，这是量子阱、量子霍尔效应的关键特征。
\end{itemize}
\end{errorbox}

% --------------------------------------------------------
% 6. 物理图像与关联模型
% --------------------------------------------------------
\begin{insightbox}[物理图像与模型关联]
\textbf{物理图像}：

金属中的自由电子就像高压水枪中的水分子：
\begin{itemize}[nosep]
    \item \textbf{费米速度} $v_F$：水分子本身以极高速度无规运动（$\sim10^6\,\text{m/s}$）。
    \item \textbf{漂移速度} $v_d$：电场相当于给水管加了一个微小的定向推力，水流的整体偏移速度很小（$\sim10^{-2}\,\text{m/s}$）。
    \item \textbf{弛豫时间} $\tau$：水分子不断与管壁（杂质、声子）碰撞，平均每隔 $10^{-14}\,\text{s}$ 改变一次方向。
\end{itemize}

\textbf{与相似模型的对比}：
\begin{center}
\begin{tabular}{llll}
\toprule
\textbf{对比维度} & \textbf{Drude} & \textbf{Sommerfeld} & \textbf{Bloch} \\
\midrule
电子统计 & 经典麦克斯韦 & 费米--狄拉克 & 费米--狄拉克 \\
能谱 & 连续 & 连续（自由电子） & 能带（周期势） \\
适用温度 & 全部 & 低温量子修正 & 完整温度范围 \\
成功之处 & 欧姆定律 & 热容、泡利顺磁性 & 完整金属理论 \\
\bottomrule
\end{tabular}
\end{center}

\textbf{推广方向}：
\begin{itemize}[nosep]
    \item 量子输运：弹道输运 ($l\gg L$)、介观电导量子化 $G=2e^2/h$
    \item 超导：Cooper 对形成后载流子变为玻色子，无电阻
    \item 拓扑物态：费米面拓扑不变量决定表面态存在性
\end{itemize}
\end{insightbox}

% --------------------------------------------------------
% 7. 推导附录
% --------------------------------------------------------
\begin{derivationbox}[推导细节（自我检验用）]
\textbf{Step 1：费米波矢推导}

$k$ 空间中费米球体积：$\frac{4}{3}\pi k_F^3$。

态密度（含自旋）：$\frac{2V}{(2\pi)^3}=\frac{V}{4\pi^3}$。

电子总数：$N=\frac{V}{4\pi^3}\cdot\frac{4}{3}\pi k_F^3=\frac{V}{3\pi^2}k_F^3$。

故 $n=N/V=\frac{k_F^3}{3\pi^2}$，即 $k_F=(3\pi^2 n)^{1/3}$。

\textbf{Step 2：Drude 电导率}

电子在电场中获得定向速度 $v_d$，碰撞后速度归零。平均自由时间内速度线性增长：
\[
\langle v_d\rangle=\frac{1}{2}\cdot\frac{-e\mathcal{E}}{m}\tau.
\]

电流密度：$j=-ne\langle v_d\rangle=\frac{ne^2\tau}{m}\mathcal{E}=\sigma\mathcal{E}$。

更严格地从玻尔兹曼方程弛豫时间近似出发，同样得到 $\sigma=ne^2\tau/m$。

\textbf{Step 3：能态密度}

3D：$k$ 空间球壳体积 $4\pi k^2dk$，态数 $dN=\frac{V}{\pi^2}k^2dk$。
由 $E=\hbar^2k^2/2m$ 得 $kdk=m dE/\hbar^2$，$k=\sqrt{2mE}/\hbar$。

\[
g_{3D}(E)=\frac{\dd N}{\dd E}=\frac{V}{\pi^2}\cdot\frac{m}{\hbar^2}\cdot\frac{\sqrt{2mE}}{\hbar}
=\frac{V}{2\pi^2}\Bigl(\frac{2m}{\hbar^2}\Bigr)^{3/2}\sqrt{E}.
\]

2D 和 1D 同理，将 $k$ 空间面积元/线元代入即可。
\end{derivationbox}

\vspace{1em}
